% from ega4-21c-fr.tex
\oldpage[IV-4]{310}
sur $k(y)$, donc fermé dans $X_y$ $(\mathbf I, 6.4.2)$. Remplaçant encore
éventuellement $X$ par un voisinage ouvert plus petit de $z$ dans $X$, on peut donc se
borner au cas où $V\cap X_y=\varnothing$, autrement dit $T\supset X_y$; comme d'autre
part $z\in X_y$ est le point générique de $T$, on a
$g(T)\subset\overline{\{y\}}=\{y\}$, autrement dit $T=X_y$. On est alors exactement
dans les conditions d'application de (21.12.10), d'où la conclusion.
\end{proof}

\begin{theorem}[21.12.12]
\label{IV.21.12.12}
\textup{(van der Waerden)}. ---
Soient $X$, $Y$ deux préschémas localement noethériens intègres, $g:X\to Y$ un
morphisme birationnel et localement de type fini. On suppose en outre que $Y$ est
normal et que pour tout $y\in Y$, l'anneau $\mathcal O_{Y,y}$ vérifie la condition (W)
de (21.12.8) (conditions remplies en particulier lorsque $Y$ est localement factoriel
(21.12.9)). Si $V$ est le plus grand ouvert de $X$ tel que la restriction
$g|V:V\to Y$ soit un isomorphisme local, toutes les composantes irréductibles de
$T=X-V$ sont de codimension $1$ dans $X$.
\end{theorem}

\begin{proof}
On notera que puisque $g$ est birationnel, l'ensemble ouvert $V$ n'est pas vide; il
suffit donc de prouver qu'un point maximal $z$ de $T$ ne peut être isolé dans $X_y$,
où $y=g(z)$. Mais, quitte à se restreindre à un voisinage ouvert de $z$, on peut se
borner au cas où $T=X_y$; alors toutes les fibres $g^{-1}(y')$ ($y'\in Y$) seraient
vides ou réduites à un point et il résulterait du « Main theorem » (8.12.10) que $g$
serait un isomorphisme local, contrairement à l'hypothèse $z\in T$.
\end{proof}

\begin{corollary}[21.12.13]
\label{IV.21.12.13}
Supposons vérifiées les hypothèses de (21.12.12) et en outre, supposons que $g$ soit
quasi-fini en chacun des points de $X^{(1)}$ (on rappelle que ce sont les points où
$\dim(\mathcal O_{X,x})=1$); alors $g$ est un isomorphisme local, et si de plus $g$ est
séparé, $g$ est une immersion ouverte.
\end{corollary}

\begin{proof}
Il suffit de prouver la première assertion, la seconde étant une conséquence de la
première et de $(\mathbf I, 8.2.8)$. Tout revient à prouver, avec les notations de
(21.12.12), que $T=\varnothing$; dans le cas contraire, un point générique $z$ d'une
composante irréductible de $T$ appartiendrait à $X^{(1)}$ en vertu de (21.12.12), et par
hypothèse il serait isolé dans $X_y$ avec $y=g(z)$; mais on a vu dans la démonstration
de (21.12.12) que cela n'est pas possible.
\end{proof}

\begin{remarks}[21.12.14]
\label{IV.21.12.14}
(i) La conclusion de (21.12.12) s'applique lorsque $Y$ est régulier et intègre et $X$
intègre, car en vertu du théorème d'Auslander--Buchsbaum (21.11.1), $Y$ est alors
localement factoriel. Par contre, on connaît des exemples où $X$ et $Y$ sont des schémas
algébriques normaux de dimension $3$, sur un corps algébriquement clos de
caractéristique quelconque, et où la conclusion de (21.12.12) n'est plus valable.

(ii) L'ensemble $T$ de (21.12.12) est aussi l'ensemble des points où $g$ est
\emph{ramifié}. En effet, si $g$ est non ramifié en un point $x\in X$, il est aussi non
ramifié dans un voisinage ouvert affine $U$ de $x$ (17.3.7), donc $g|U$ est un morphisme
séparé, quasi-fini et birationnel (17.4.3). Puisque $Y$ est normal, on conclut du « Main
theorem » $(\mathbf{III}, 4.4.9)$ que $g$ est un isomorphisme local au point $x$, donc
$x\in X-T$. Réciproquement, $g$ est évidemment non ramifié en tout point où il est un
isomorphisme local. Ceci justifie le titre donné à cette section.

